本部分讨论智能体的「大脑」，以及它们如何用大脑来学习。  
上一部分展示了基础的下一个词元预测式大语言模型是如何训练的。在智能体模型中，通常会使用更高级的推理模型与多模态模型（我们仍统称为{\em 大语言模型}）。本部分将讨论这些更高级的模型、为什么需要它们，以及它们如何被训练与使用。  

本部分第一章介绍推理模型。  
最初大语言模型无法解决简单的推理问题，而思维链这一概念的提出，让该领域迈出了一大步。我们将讨论数学应用题、思维链提示与上下文学习。我们将看到一些模型利用自我反思来尝试改进自身行为。自我反思循环催生了面向大语言模型的智能体运行框架。  

在随后一章中，我们走出语言：我们将看到 Transformer 架构同样适用于视觉与动作领域。我们将讨论各类世界模型——它们让智能体在与世界交互时能更智能地行动。  

本部分最后一章介绍元学习这一概念，它是智能体大语言模型中许多学习类型共有的范式。我们将讨论什么是元学习，以及为什么它是一种颇具挑战性的机器学习形式。
